\section{Methodology}
\label{sec:method}

Building on the insight that oracle semantics must be decoupled from input strategy, we design a dual-experiment framework.

\subsection{Overview}

Our methodology separates two sources of potential oracle advantage:

\begin{enumerate}
\item \textbf{Oracle semantic strength} --- the contract oracle detects violations that the differential oracle \emph{cannot} by construction, regardless of inputs used (measured by Experiment A).
\item \textbf{Input exploration advantage} --- adaptive PBT explores more of the input space than static test sets (measured by Experiment B).
\end{enumerate}

\subsection{Dataset: ContractEval}

ContractEval \citep{Lim2025ContractEval} provides 364 tasks from HumanEval+ and MBPP+ with Python inline \texttt{assert} pre/post-conditions. Each task includes CVT inputs ($\sim$5 per task) specifically designed to trigger contract violations. Zero tasks were quarantined by our soundness pre-check.

\subsection{LLM Corpus}

Five model families: GPT-4o-mini, Claude-3-Haiku (closed); DeepSeek-Coder-V2-Lite (16B), CodeLlama-13B, CodeLlama-34B (open). Code samples ($n=10$ per model per task) from the evalplus generation corpus, filtered to test-passing programs. Experiment A (oracle isolation) completed for 3 of 5 model families (Claude-3-Haiku, CodeLlama-13B, CodeLlama-34B; 10,432 triples); Experiment B (adaptive PBT) completed for all 5 families (17,226 triples).

\subsection{Experiment A: Oracle Isolation via CVT Inputs}

ContractEval's inline \texttt{assert} contracts are precondition checks that fire only on invalid inputs. EvalPlus's 764 static inputs are valid by construction --- producing zero contract violations. We use CVT inputs as the controlled comparison set.

\begin{itemize}
\item \textbf{Oracle A (differential):} $p(x) \neq \mathrm{gt\_plain}(x)$ --- did LLM output differ from plain canonical?
\item \textbf{Oracle B (contract):} \texttt{canonical\_with\_contract(x)} raises \texttt{AssertionError} --- did the contract fire?
\item \textbf{Contract-unique (CU):} Oracle A = PASS \emph{and} Oracle B = FAIL. Programs that output the ``correct'' answer while violating the contract invariant.
\item \textbf{Oracle-isolation gap} = mean(Oracle B rate) $-$ mean(Oracle A rate) over all triples.
\end{itemize}

Implementation: 8-process multiprocessing pool; 5s per-input timeout; 10s per-program timeout; JSONL streaming.

\subsection{Experiment B: Adaptive PBT Contribution}

icontract-hypothesis (5,000 examples per triple; 30s SIGALRM timeout) explores contract-relevant input space beyond CVT coverage. Adaptive contribution per triple = PBT failure rate $-$ static CVT failure rate. 94.6\% compatibility rate (17,226/18,200 triples valid).

\subsection{Contract Richness Analysis (Experiment C)}

4-tier AST-based richness scoring: Tier 1 (Simple), Tier 2 (BoolOp chaining), Tier 3 (any/all quantification), Tier 4 (Compound). Spearman $\rho$ correlation with oracle-isolation gap; bootstrap CI; permutation test.

\subsection{Cross-Model Analysis (Experiment D)}

Kendall $\tau$ (contract-satisfaction ranking vs.\ pass@1*); MixedLM partial $\Delta R^2$ for model identity; exact permutation test (type=pairings). Statistical note: at $n=5$, minimum achievable $p \approx 0.017$ for two-sided test at $|\tau|=1.0$ only.
